(1) 在惯性系对火箭加尾气的小时间控制体作动量平衡，排气携带相对速度 \(-u\)，因此 \(m\,dv/dt=u\alpha-mg\)，\(dm/dt=-\alpha\)。由 \(m(t)=m_0-\alpha t\) 积分得 \(v(t)=u\ln[m_0/m(t)]-gt\)，加速度为 \(a(t)=u\alpha/m(t)-g\)。对数项不是经验修正，而是质量变动与给定相对排气速度共同导致。
(2) 因为 \(\int_0^t\ln[m_0/(m_0-\alpha s)]\,ds=t-[m(t)/\alpha]\ln[m_0/m(t)]\)，故 \(z(t)=u\{t-[m(t)/\alpha]\ln[m_0/m(t)]\}-gt^2/2\)，并满足初值 \(z(0)=\dot z(0)=0\)。将该式直接求导可恢复前问的速度，从而检查了积分常数与括号中的质量量纲。
(3) 发射架释放时若必须从静止立即离地，要求初始推力 \(u\alpha>m_0g\)；否则在解除约束前不能把数学上的负高度解释为真实飞行。停火时 \(v_b=u\ln(m_0/m_f)-gt_b\)，\(z_b=u[t_b-m_f\ln(m_0/m_f)/\alpha]-gt_b^2/2\)。如果 \(v_b>0\)，其后的弹道峰值为 \(z_{\max}=z_b+v_b^2/(2g)\)；若 \(v_b\le0\) 则停火瞬间已不再上升。
(4) 反作用力的模是 \(F_{\rm thrust}=u\alpha\)，确实恒定，但惯性质量 \(m(t)\) 逐渐减小，故 \(a=F_{\rm thrust}/m(t)-g\) 不是常数，直接写 \(v=at\) 是错误的。对数律需尾气相对喷口速度恒定、喷流稳态且没有压力推力修正；空气阻力、随高度变重力、燃料分级会修改方程。
